% Copy into a results folder; replace every placeholder before claiming completion.
% Paths below are relative to the repository root.
% id: R-prefix-001
% title: [Descriptive mathematical title]
% stage: ideas
% proof-status: direction-only
% Other statuses: conjectural | partial | complete-unreviewed | reviewed | refuted
% depends-on: [result IDs, definition labels, and source keys; or none]
% derived-from: [origin result IDs; or none]
% computation: [code paths and reproduction notes; or none]
% human-review: not reviewed
% Review record when applicable: name | YYYY-MM-DD | reviewed statement/proof version
% Append contributions; never overwrite earlier attribution:
% YYYY-MM-DD | human/operator: [name] | service/app: [name or none]
% model: [reported identifier, unknown, or none for human-only work]
% role: [originated / drafted / revised / checked] | scope: [specific contribution]
% Optional: session reference or commit; never invent unavailable identifiers.

\section{Descriptive mathematical title}\label{sec:prefix-001}

\paragraph{Setting and purpose.}
\emph{Specify the objects, hypotheses, relevant definitions, dependencies,
and the role of the proposed result in the project.}

% Select question/conjecture/lemma/proposition/theorem as appropriate.
% Change the environment as the result develops; keep the label stable.
\begin{conjecture}[Working statement]\label{res:prefix-001}
\emph{Replace with a precise mathematical statement, including quantifiers
and every additional hypothesis. This placeholder asserts no result.}
\end{conjecture}

% Early stages: direction or sketch, with no end-of-proof symbol.
% Once the proof is complete, replace this paragraph with a proof environment.
\paragraph{Proof direction (incomplete).}
\emph{Describe the proposed argument and its dependencies. Number nontrivial
steps when this makes gaps easier to reference.}

\paragraph{Gaps and checks.}
\emph{List unresolved steps, edge cases, citation checks, and the limits of
computational evidence. Do not erase gaps merely to advance a stage.}

\paragraph{Consequences.}
\emph{Explain what follows if the claim is established, and its relevance
to the main research question.}

% For paper-ready exposition, resolve placeholders and gaps, supply a complete
% proof where one is claimed, and adapt/remove these paragraph headings.
% Retain provenance in comments. Substantive edits invalidate earlier review.
